\chapter{Complexity — Pair Encode}
%%% generated by the blueprint pipeline from TCSlib.Complexity.CircuitComplexity.PairEncode
\section{Overview}

Declarations of the Lean module \texttt{TCSlib.Complexity.CircuitComplexity.PairEncode}.

\section{Declarations}

\begin{definition}[Input and constant gate construction]
\lean{BoolCircuit.inputGate}
\label{BoolCircuit.inputGate}
\leanok
\uses{BoolCircuit.DAGGate, BoolCircuit.constGate}
For a natural number $n$, define a gate from either an input index $i\in\mathrm{Fin}(n)$
or a Boolean value $b$. An input index gives an AND gate whose sole argument is the
natural-number value of $i$, while $b$ gives the constant gate with value $b$.
\end{definition}

\begin{definition}[Input buffer gates]
\lean{BoolCircuit.inputBuffer}
\label{BoolCircuit.inputBuffer}
\leanok
\uses{BoolCircuit.DAGGate, BoolCircuit.inputGate}
For natural numbers $n,m$ and a map $s$ assigning to each index $i<m$ either an input
index below $n$ or a Boolean constant, the input buffer is the list of $m$ gates
obtained by applying the corresponding input-or-constant gate to each $s(i)$ in order.
\end{definition}

\begin{definition}[Boolean input values]
\lean{BoolCircuit.inputValues}
\label{BoolCircuit.inputValues}
\leanok
Let $n,m$ be nonnegative integers, let $s$ assign to each of $m$ input positions either
an index in $\text{Fin}(n)$ or a Boolean constant, and let $x:\text{Fin}(n)\to\text{Bool}$
be an input assignment. The resulting assignment on the $m$ positions takes the value
$x(j)$ when $s$ selects index $j$, and the selected constant otherwise.
\end{definition}

\begin{theorem}[Input Gate Argument Bounds]
\lean{BoolCircuit.inputGate_args_lt}
\label{BoolCircuit.inputGate_args_lt}
\leanok
\uses{BoolCircuit.inputGate}
\difficulty{4}
Let $n$ be a natural number and let a source selector be either an input index among $n$
inputs or a Boolean constant. Every argument index of the gate associated with that
selector is less than $n$.
\end{theorem}

\begin{definition}[Buffering circuit inputs]
\lean{BoolCircuit.DAGCircuit.bufferInputs}
\label{BoolCircuit.DAGCircuit.bufferInputs}
\leanok
\uses{BoolCircuit.DAGCircuit, BoolCircuit.DAGGate.remap, BoolCircuit.inputBuffer, BoolCircuit.inputGate_args_lt}
Let $n,m$ be nonnegative integers, let $C$ be a directed acyclic Boolean circuit with
$m$ input vertices, and let $s$ assign to each of those inputs either one of $n$ new
input vertices or a Boolean constant. The circuit on $n$ inputs formed by first adding
input gates for the assignment specified by $s$, then appending the gates of $C$ with
all their referenced vertices shifted by $n$, has output vertex $n$ plus the output
vertex of $C$.
\end{definition}

\begin{theorem}[Values Produced by an Input Buffer]
\lean{BoolCircuit.inputBuffer_value}
\label{BoolCircuit.inputBuffer_value}
\leanok
\uses{BoolCircuit.DAGGate.eval, BoolCircuit.inputBuffer, BoolCircuit.inputGate, BoolCircuit.inputValues, BoolCircuit.runWith, BoolCircuit.runWith_getD_gate, BoolCircuit.runWith_getD_of_lt}
\difficulty{5}
Let $n,m$ be natural numbers, let each of $m$ source positions select either one of $n$
Boolean inputs or a Boolean constant, and let $x$ be an assignment of Boolean values to
those $n$ inputs. Form the list of gates that copies the selected input or produces the
selected constant, and run these gates in order starting with the list of the $n$ input
values, evaluating each gate from the values already present. For every buffer position
$i$, the value at index $n+i$ of the resulting list, using false as the default if out
of range, equals the value selected by source position $i$.
\end{theorem}

\begin{theorem}[Evaluation after buffering inputs]
\lean{BoolCircuit.DAGCircuit.bufferInputs_eval}
\label{BoolCircuit.DAGCircuit.bufferInputs_eval}
\leanok
\uses{BoolCircuit.DAGCircuit, BoolCircuit.DAGCircuit.bufferInputs, BoolCircuit.DAGCircuit.eval, BoolCircuit.DAGCircuit.values, BoolCircuit.DAGGate.eval, BoolCircuit.DAGGate.eval_remap, BoolCircuit.inputBuffer, BoolCircuit.inputBuffer_value, BoolCircuit.inputValues, BoolCircuit.runWith, BoolCircuit.runWith_append, BoolCircuit.runWith_remap_rel}
\difficulty{4}
Let $C$ be a directed acyclic Boolean circuit with $m$ inputs, let $s$ assign to each of
its inputs either one of $n$ input vertices or a fixed Boolean constant, and let $x$ be
an assignment of Boolean values to those $n$ inputs. The circuit formed by buffering the
inputs according to $s$ evaluates on $x$ to the same value as $C$ evaluates on the
assignment obtained by substituting the selected input values and constants.
\end{theorem}

\begin{theorem}[Structure of the buffered circuit]
\lean{BoolCircuit.DAGCircuit.bufferInputs_structure}
\label{BoolCircuit.DAGCircuit.bufferInputs_structure}
\leanok
\uses{BoolCircuit.DAGCircuit, BoolCircuit.DAGCircuit.IsFaninTwo, BoolCircuit.DAGCircuit.IsWellFormed, BoolCircuit.DAGCircuit.bufferInputs, BoolCircuit.DAGCircuit.size, BoolCircuit.DAGGate.FaninTwo, BoolCircuit.DAGGate.WellFormed, BoolCircuit.DAGGate.remap, BoolCircuit.inputBuffer, BoolCircuit.inputGate}
\difficulty{5}
Let $C$ be a directed acyclic Boolean circuit with $m$ inputs, let $s$ assign each of
its inputs either one of $n$ input vertices or a Boolean constant, and suppose every
gate of $C$ has pairwise distinct arguments, every negation gate has exactly one
argument, and every gate has at most two arguments. Then the buffered circuit has
pairwise distinct arguments at each gate, every negation gate has exactly one argument,
every gate has at most two arguments, and its size is $n+\mathrm{size}(C)$.
\end{theorem}

\begin{definition}[Paired wire list]
\lean{BoolCircuit.pairWireList}
\label{BoolCircuit.pairWireList}
\leanok
For a nonnegative integer $n$ and a finite Boolean list $r$, form a list of wires by
listing each of the $n$ input indices twice, in order, and then listing the Boolean
constants false and true followed by the entries of $r$, each as a constant wire.
\end{definition}

\begin{theorem}[Length and evaluation of the pair wire list]
\lean{BoolCircuit.pairWireList_spec}
\label{BoolCircuit.pairWireList_spec}
\leanok
\uses{BoolCircuit.pairWireList, Turing.length_pairEncode, Turing.pairEncode}
\difficulty{4}
For every natural number $n$ and bit string $r$, the pair wire list has length
$2n+2+|r|$. Moreover, for every assignment $v:\text{Fin}(n)\to\text{Bool}$, replacing each
variable wire by its value and each constant wire by that constant gives the
self-delimiting pairing of the bit string $(v(0),\ldots,v(n-1))$ with $r$.
\end{theorem}

\begin{definition}[Function represented by a list]
\lean{BoolCircuit.vectorOfList}
\label{BoolCircuit.vectorOfList}
\leanok
Let $\alpha$ be a type, let $m$ be a nonnegative integer, and let $l$ be a list of
elements of $\alpha$ of length $m$. The function on $\text{Fin}(m)$ represented by $l$
returns its entry at the corresponding position.
\end{definition}

\begin{theorem}[Reconstruction of a List from Its Vector]
\lean{BoolCircuit.vectorOfList_ofFn}
\label{BoolCircuit.vectorOfList_ofFn}
\leanok
\uses{BoolCircuit.vectorOfList}
\difficulty{3}
Let $m$ be a natural number, let $l$ be a list of elements of a type, and suppose its
length is $m$. Form the function from indices $0,\ldots,m$ to elements of the type by reading
the corresponding entries of $l$; applying the list constructor indexed by that function
recovers $l$.
\end{theorem}

\begin{definition}[Paired wire assignment]
\lean{BoolCircuit.pairWiring}
\label{BoolCircuit.pairWiring}
\leanok
\uses{BoolCircuit.pairWireList, BoolCircuit.pairWireList_spec, BoolCircuit.vectorOfList}
For a nonnegative integer $n$ and a finite Boolean list $r$, the paired wire assignment
on $2n+2+|r|$ positions is the function whose list of values is obtained by listing
every input index twice in order, then the constants false and true, then the entries of
$r$ as constant wires.
\end{definition}

\begin{definition}[Paired input encoding]
\lean{BoolCircuit.pairEncodeInput}
\label{BoolCircuit.pairEncodeInput}
\leanok
\uses{BoolCircuit.inputValues, BoolCircuit.pairWiring}
Let $n$ be a nonnegative integer, let $r$ be a finite Boolean list, and let $v$ be a
Boolean assignment on the $n$ input indices. The paired input encoding on $2n+2+|r|$
positions evaluates the paired wire assignment using $v$ for input wires and the
specified values for constant wires.
\end{definition}

\begin{definition}[Paired circuit encoding]
\lean{BoolCircuit.DAGCircuit.pairEncode}
\label{BoolCircuit.DAGCircuit.pairEncode}
\leanok
\uses{BoolCircuit.DAGCircuit, BoolCircuit.DAGCircuit.bufferInputs, BoolCircuit.pairWiring}
Let $n$ be a nonnegative integer, let $r$ be a finite Boolean list, and let $C$ be a
directed acyclic Boolean circuit with $2n+2+|r|$ inputs, gates in topological order,
valid earlier-vertex references, and a valid output vertex. Its paired encoding is the
circuit on $n$ inputs obtained by buffering its inputs according to the paired wire
assignment.
\end{definition}

\begin{theorem}[Input vector for paired encoding]
\lean{BoolCircuit.pairEncodeInput_ofFn}
\label{BoolCircuit.pairEncodeInput_ofFn}
\leanok
\uses{BoolCircuit.pairEncodeInput, BoolCircuit.pairWireList_spec, BoolCircuit.pairWiring, BoolCircuit.vectorOfList_ofFn, Turing.pairEncode}
\difficulty{3}
Let $n$ be a natural number, let $r$ be a finite bit string, and let $v$ assign a bit to
each of the $n$ input positions. The finite vector of bits obtained by evaluating the
wiring for the paired encoding of $v$ and $r$ is exactly the paired encoding of the list
of values of $v$ and the string $r$.
\end{theorem}

\begin{theorem}[Evaluation under pair encoding]
\lean{BoolCircuit.DAGCircuit.pairEncode_eval}
\label{BoolCircuit.DAGCircuit.pairEncode_eval}
\leanok
\uses{BoolCircuit.DAGCircuit, BoolCircuit.DAGCircuit.bufferInputs_eval, BoolCircuit.DAGCircuit.eval, BoolCircuit.DAGCircuit.pairEncode, BoolCircuit.pairEncodeInput, BoolCircuit.pairWiring}
\difficulty{1}
Let $n$ be a nonnegative integer, let $r$ be a finite bit string, and let $C$ be a
Boolean circuit with $2n+2+|r|$ inputs, whose gates read only earlier vertices and whose
output is a vertex of the circuit. For every assignment $v$ of Boolean values to the $n$
inputs, the circuit obtained by wiring each input bit twice, followed by $01$ and then
$r$, evaluates to the same Boolean value as $C$ evaluated on that encoded input.
\end{theorem}

\begin{theorem}[Pair Encoding Preserves Fan-In Two]
\lean{BoolCircuit.DAGCircuit.pairEncode_isFaninTwo}
\label{BoolCircuit.DAGCircuit.pairEncode_isFaninTwo}
\leanok
\uses{BoolCircuit.DAGCircuit, BoolCircuit.DAGCircuit.IsFaninTwo, BoolCircuit.DAGCircuit.bufferInputs_structure, BoolCircuit.DAGCircuit.pairEncode, BoolCircuit.pairWiring}
\difficulty{1}
Let $n$ be a nonnegative integer, let $r$ be a finite bit string, and let $C$ be a
Boolean circuit with $2n+2+|r|$ inputs and gates in topological order, each reading only
earlier vertices, with a valid output vertex. Suppose every gate of $C$ has distinct
input vertices, every NOT gate has exactly one input, and every gate has at most two
inputs. Then the circuit obtained by wiring its inputs to the pair encoding of an
$n$-bit input and $r$ also has distinct inputs at each gate, exactly one input at every
NOT gate, and at most two inputs at every gate.
\end{theorem}

\begin{theorem}[Pair Encoding Preserves Well-Formedness]
\lean{BoolCircuit.DAGCircuit.pairEncode_isWellFormed}
\label{BoolCircuit.DAGCircuit.pairEncode_isWellFormed}
\leanok
\uses{BoolCircuit.DAGCircuit, BoolCircuit.DAGCircuit.IsFaninTwo, BoolCircuit.DAGCircuit.IsWellFormed, BoolCircuit.DAGCircuit.bufferInputs_structure, BoolCircuit.DAGCircuit.pairEncode, BoolCircuit.pairWiring}
\difficulty{1}
Let $n$ be a nonnegative integer, let $r$ be a finite bit string, and let $C$ be a
Boolean circuit with $2n+2+|r|$ inputs, gates in topological order, and a valid output
vertex. Suppose every gate has distinct input vertices, every NOT gate has exactly one
input, and every gate has at most two inputs. Then the circuit obtained by wiring its
inputs to the pair encoding of an $n$-bit input and $r$ has distinct input vertices at
each gate, and every NOT gate has exactly one input.
\end{theorem}

\begin{theorem}[Size After Pair Encoding]
\lean{BoolCircuit.DAGCircuit.pairEncode_size}
\label{BoolCircuit.DAGCircuit.pairEncode_size}
\leanok
\uses{BoolCircuit.DAGCircuit, BoolCircuit.DAGCircuit.bufferInputs, BoolCircuit.DAGCircuit.pairEncode, BoolCircuit.DAGCircuit.size, BoolCircuit.inputBuffer}
\difficulty{1}
Let $n$ be a nonnegative integer, let $r$ be a finite bit string, and let $C$ be a
Boolean circuit with $2n+2+|r|$ inputs, gates in topological order, and a valid output
vertex. If circuit size is the number of inputs plus the number of gates, then wiring
$C$ to the pair encoding of an $n$-bit input and $r$ produces a circuit of size
$n+\operatorname{size}(C)$.
\end{theorem}

\begin{theorem}[Acceptance of a finite assignment]
\lean{BoolCircuit.DAGCircuitFamily.mem_language_ofFn}
\label{BoolCircuit.DAGCircuitFamily.mem_language_ofFn}
\leanok
\uses{BoolCircuit.DAGCircuit.eval, BoolCircuit.DAGCircuitFamily, BoolCircuit.DAGCircuitFamily.language, BoolCircuit.DAGCircuitFamily.mem_language_iff}
\difficulty{3}
Let $F$ be a family of Boolean circuits, let $m$ be a natural number, and let $v$ be a
Boolean assignment to the $m$ input positions. The length-$m$ word whose bits are given
by $v$ belongs to the language accepted by $F$ if and only if the circuit in $F$ for
inputs of length $m$ outputs true on $v$.
\end{theorem}

\begin{theorem}[Evaluation on paired inputs]
\lean{BoolCircuit.DAGCircuitFamily.pairEncode_eval_eq_true_iff}
\label{BoolCircuit.DAGCircuitFamily.pairEncode_eval_eq_true_iff}
\leanok
\uses{BoolCircuit.DAGCircuit.eval, BoolCircuit.DAGCircuitFamily, BoolCircuit.DAGCircuitFamily.language, BoolCircuit.DAGCircuitFamily.mem_language_ofFn, BoolCircuit.pairEncodeInput, BoolCircuit.pairEncodeInput_ofFn, Turing.pairEncode}
\difficulty{3}
Let $F$ be a family of Boolean circuits, let $n$ be a natural number, let $r$ be a bit
string, and let $v$ assign a bit to each of the $n$ input positions. The circuit in $F$
for inputs of length $2n+2+|r|$, evaluated on the input formed by duplicating each bit
of $v$, then appending $01$ and $r$, outputs $1$ if and only if the paired encoding of
the bit string specified by $v$ and $r$ belongs to the language accepted by $F$.
\end{theorem}
